% Anexo — Contratos de la API REST
% Documenta los endpoints principales agrupados por módulo funcional.

Este anexo presenta los contratos de los endpoints principales de la API REST del sistema DOMU, organizados por módulo funcional. Para cada módulo se incluye una tabla con los métodos HTTP, rutas, descripción y los \textit{Data Transfer Objects} (DTOs) de solicitud y respuesta, seguida de la definición de los campos principales de cada DTO en formato de clase Java.

% ============================================================
\section{Módulo A --- Visitas}
% ============================================================

El módulo de visitas gestiona el registro de ingreso y salida de visitantes, la verificación de pre-autorizaciones mediante código QR y la creación rápida de visitas desde contactos frecuentes.

\begin{table}[H]
\centering
\small
\begin{tabular}{|p{0.08\textwidth}|p{0.35\textwidth}|p{0.25\textwidth}|p{0.12\textwidth}|p{0.12\textwidth}|}
\hline
\textbf{Método} & \textbf{Ruta} & \textbf{Descripción} & \textbf{Request} & \textbf{Response} \\\hline
POST & \texttt{/api/visits} & Registrar nueva visita & CreateVisit\-Request & Visit\-Response \\\hline
POST & \texttt{/api/visits/qr-check} & Verificar QR de pre-autorización & QrCheck\-Request & QrCheck\-Response \\\hline
POST & \texttt{/api/visits/\{id\}/check-in} & Registrar check-in de visita & --- & Visit\-Response \\\hline
POST & \texttt{/api/visit-contacts/\{id\}/register} & Crear visita desde contacto frecuente & --- & Visit\-Response \\\hline
\end{tabular}
\caption{Endpoints del módulo de Visitas.}
\label{tab:endpoints-visitas}
\end{table}

\begin{lstlisting}[style=javadto, title={\texttt{CreateVisitRequest}}]
public class CreateVisitRequest {
    private String visitorName;           // Nombre completo del visitante
    private String visitorRut;            // RUT o pasaporte del visitante
    private Long unitId;                  // ID de la unidad habitacional destino
    private String reason;                // Motivo de la visita
    private LocalDateTime scheduledDate;  // Fecha programada (opcional, pre-autorizacion)
}
\end{lstlisting}

\begin{lstlisting}[style=javadto, title={\texttt{VisitResponse}}]
public class VisitResponse {
    private Long id;                      // Identificador de la visita
    private String visitorName;           // Nombre completo del visitante
    private String visitorRut;            // RUT o pasaporte del visitante
    private String status;                // Estado: PENDING, CHECKED_IN, CHECKED_OUT
    private LocalDateTime checkInTime;    // Marca temporal de ingreso
    private LocalDateTime checkOutTime;   // Marca temporal de salida
    private String unitNumber;            // Numero de la unidad destino
}
\end{lstlisting}

\begin{lstlisting}[style=javadto, title={\texttt{QrCheckRequest}}]
public class QrCheckRequest {
    private String qrPayload;  // Cadena leida del codigo QR de la cedula
}
\end{lstlisting}

\begin{lstlisting}[style=javadto, title={\texttt{QrCheckResponse}}]
public class QrCheckResponse {
    private boolean authorized;  // Indica si la pre-autorizacion es vigente
    private Long visitId;        // ID de la visita asociada (si existe)
    private String visitorName;  // Nombre del visitante pre-autorizado
    private String visitorRut;   // RUT del visitante
}
\end{lstlisting}

% === PLACEHOLDER: Captura de pantalla Visitas ===
\begin{figure}[H]
  \centering
  \fbox{\parbox{0.85\textwidth}{\centering\vspace{2.5cm}
  \textit{[Insertar captura de pantalla: Vista del conserje --- panel de registro de visitas con lectura QR]}
  \vspace{2.5cm}}}
  \caption{Vista del módulo de Visitas en el Frontend.}
  \label{fig:anexo-visitas-frontend}
\end{figure}

% ============================================================
\section{Módulo B --- Encomiendas}
% ============================================================

El módulo de encomiendas permite al conserje registrar paquetes con evidencia fotográfica, actualizar su estado y al residente consultar sus encomiendas pendientes.

\begin{table}[H]
\centering
\small
\begin{tabular}{|p{0.08\textwidth}|p{0.35\textwidth}|p{0.25\textwidth}|p{0.12\textwidth}|p{0.12\textwidth}|}
\hline
\textbf{Método} & \textbf{Ruta} & \textbf{Descripción} & \textbf{Request} & \textbf{Response} \\\hline
POST & \texttt{/api/parcels} & Registrar nueva encomienda & CreateParcel\-Request & Parcel\-Response \\\hline
PUT & \texttt{/api/parcels/\{id\}/status} & Actualizar estado de encomienda & UpdateStatus\-Request & Parcel\-Response \\\hline
GET & \texttt{/api/parcels/my} & Listar encomiendas del residente & --- & List<Parcel\-Response> \\\hline
\end{tabular}
\caption{Endpoints del módulo de Encomiendas.}
\label{tab:endpoints-encomiendas}
\end{table}

\begin{lstlisting}[style=javadto, title={\texttt{CreateParcelRequest}}]
public class CreateParcelRequest {
    private Long unitId;            // Unidad habitacional destinataria
    private String description;     // Descripcion del paquete
    private String senderName;      // Nombre del remitente o empresa de envio
    private MultipartFile photo;    // Fotografia de evidencia (opcional)
}
\end{lstlisting}

\begin{lstlisting}[style=javadto, title={\texttt{ParcelResponse}}]
public class ParcelResponse {
    private Long id;                       // Identificador de la encomienda
    private String description;            // Descripcion del paquete
    private String senderName;             // Nombre del remitente
    private String status;                 // Estado: RECEIVED, NOTIFIED, DELIVERED
    private String photoUrl;               // URL de la foto almacenada en Box
    private LocalDateTime receivedAt;      // Marca temporal de recepcion
    private LocalDateTime deliveredAt;     // Marca temporal de entrega
}
\end{lstlisting}

\begin{lstlisting}[style=javadto, title={\texttt{UpdateStatusRequest}}]
public class UpdateStatusRequest {
    private String status;  // Nuevo estado de la encomienda
}
\end{lstlisting}

% === PLACEHOLDER: Captura de pantalla Encomiendas ===
\begin{figure}[H]
  \centering
  \fbox{\parbox{0.85\textwidth}{\centering\vspace{2.5cm}
  \textit{[Insertar captura de pantalla: Vista del conserje --- registro de encomienda con foto de evidencia]}
  \vspace{2.5cm}}}
  \caption{Vista del módulo de Encomiendas en el Frontend.}
  \label{fig:anexo-encomiendas-frontend}
\end{figure}

% ============================================================
\section{Módulo C --- Gastos Comunes}
% ============================================================

El módulo financiero gestiona la creación de períodos de gasto común, el procesamiento de pagos, la generación de documentos PDF y la subida de comprobantes de transferencia.

\begin{table}[H]
\centering
\small
\begin{tabular}{|p{0.08\textwidth}|p{0.35\textwidth}|p{0.25\textwidth}|p{0.12\textwidth}|p{0.12\textwidth}|}
\hline
\textbf{Método} & \textbf{Ruta} & \textbf{Descripción} & \textbf{Request} & \textbf{Response} \\\hline
POST & \texttt{/api/common-expenses/periods} & Crear período de gasto común & CreatePeriod\-Request & Period\-Response \\\hline
POST & \texttt{/api/payments/pay-simulated} & Procesar pago simulado & PaySimulated\-Request & Payment\-Response \\\hline
GET & \texttt{/api/common-expenses/\{id\}/pdf} & Descargar estado de cuenta PDF & --- & Archivo PDF \\\hline
POST & \texttt{/api/payments/\{id\}/receipt} & Subir comprobante de transferencia & Multipart\-File & Payment\-Response \\\hline
\end{tabular}
\caption{Endpoints del módulo de Gastos Comunes.}
\label{tab:endpoints-gastos}
\end{table}

\begin{lstlisting}[style=javadto, title={\texttt{CreatePeriodRequest}}]
public class CreatePeriodRequest {
    private int month;                  // Mes del periodo (1--12)
    private int year;                   // Ano del periodo
    private List<ChargeItem> charges;   // Lista de cargos por unidad
    // ChargeItem: unitId, amount, description
}
\end{lstlisting}

\begin{lstlisting}[style=javadto, title={\texttt{PeriodResponse}}]
public class PeriodResponse {
    private Long id;                          // Identificador del periodo
    private int month;                        // Mes del periodo
    private int year;                         // Ano del periodo
    private BigDecimal totalAmount;           // Monto total del periodo
    private List<ChargeResponse> charges;     // Detalle de cargos por unidad
}
\end{lstlisting}

\begin{lstlisting}[style=javadto, title={\texttt{PaySimulatedRequest}}]
public class PaySimulatedRequest {
    private List<Long> chargeIds;     // IDs de los cargos a pagar (pago parcial)
    private String paymentMethod;     // Metodo de pago: CARD, TRANSFER
}
\end{lstlisting}

\begin{lstlisting}[style=javadto, title={\texttt{PaymentResponse}}]
public class PaymentResponse {
    private Long id;                    // Identificador del pago
    private BigDecimal amount;          // Monto pagado
    private String status;              // Estado: PENDING, COMPLETED
    private String receiptUrl;          // URL del comprobante almacenado
    private LocalDateTime paidAt;       // Marca temporal del pago
}
\end{lstlisting}

% === PLACEHOLDER: Captura de pantalla Gastos Comunes ===
\begin{figure}[H]
  \centering
  \fbox{\parbox{0.85\textwidth}{\centering\vspace{2.5cm}
  \textit{[Insertar captura de pantalla: Vista del residente --- flujo de pago de gastos comunes]}
  \vspace{2.5cm}}}
  \caption{Vista del módulo de Gastos Comunes en el Frontend.}
  \label{fig:anexo-gastos-frontend}
\end{figure}

% ============================================================
\section{Módulo D --- Residentes y Comunidad}
% ============================================================

El módulo de residentes y comunidad cubre la invitación de administradores y la consulta del directorio de residentes vinculados a una comunidad.

\begin{table}[H]
\centering
\small
\begin{tabular}{|p{0.08\textwidth}|p{0.35\textwidth}|p{0.25\textwidth}|p{0.12\textwidth}|p{0.12\textwidth}|}
\hline
\textbf{Método} & \textbf{Ruta} & \textbf{Descripción} & \textbf{Request} & \textbf{Response} \\\hline
POST & \texttt{/api/buildings/invite-admin} & Invitar administrador a comunidad & InviteAdmin\-Request & Message\-Response \\\hline
GET & \texttt{/api/buildings/\{id\}/residents} & Listar residentes de la comunidad & --- & List<Resident\-Response> \\\hline
\end{tabular}
\caption{Endpoints del módulo de Residentes y Comunidad.}
\label{tab:endpoints-residentes}
\end{table}

\begin{lstlisting}[style=javadto, title={\texttt{InviteAdminRequest}}]
public class InviteAdminRequest {
    private String email;       // Correo electronico del administrador a invitar
    private Long buildingId;    // Identificador de la comunidad
}
\end{lstlisting}

\begin{lstlisting}[style=javadto, title={\texttt{ResidentResponse}}]
public class ResidentResponse {
    private Long id;              // Identificador del usuario
    private String fullName;      // Nombre completo del residente
    private String email;         // Correo electronico
    private String unitNumber;    // Numero de la unidad habitacional asignada
    private String role;          // Rol del usuario dentro de la comunidad
}
\end{lstlisting}

% === PLACEHOLDER: Captura de pantalla Residentes ===
\begin{figure}[H]
  \centering
  \fbox{\parbox{0.85\textwidth}{\centering\vspace{2.5cm}
  \textit{[Insertar captura de pantalla: Vista del administrador --- listado de residentes de la comunidad]}
  \vspace{2.5cm}}}
  \caption{Vista del módulo de Residentes y Comunidad en el Frontend.}
  \label{fig:anexo-residentes-frontend}
\end{figure}
